% Preambuła wspólna dla fragmentów testowych — parametry z gen_book.md §7.
\usepackage{lmodern}
\usepackage{tikz}
\usetikzlibrary{arrows.meta}
\usepackage{booktabs}

% Wcięcie akapitu zamiast odstępu — §3.1 pkt 1.
\setlength{\parindent}{1em}
\setlength{\parskip}{0pt}

% Bez żywej paginy: numer strony na dole, tytuł rozdziału nigdzie.
%
% Klasa `book` domyślnie stawia nad tekstem pasek z tytułem rozdziału i numerem
% strony. Przy marginesie górnym 18 pt pasek **nie mieści się nad kolumną
% i ląduje poza papierem** — w PDF-ie siedzi na współrzędnej y = -8 pt, ucięty.
% Nie widać tego w tekście wyciągniętym z pliku ani w logu kompilacji; widać
% dopiero, gdy ktoś zauważy, że dokument nie ma numerów stron.
%
% Można to naprawić opcją `includehead` geometry, ale na ekranie Kindle'a
% mieszczącym 24 wiersze żywa pagina zjada wiersz, nie wnosząc nic: tytuł
% rozdziału czytelnik ma w spisie treści, a pozycję pokazuje czytnik.
\pagestyle{plain}

% Wąska kolumna wymaga tolerancyjnego łamania: przy 53 znakach w wierszu
% domyślne ustawienia TeX-a zostawiają rzeki białych plam.
\tolerance=1500
\emergencystretch=1em

% Strona może kończyć się wyżej, zamiast rozciągać odstępy.
%
% Domyślne `\flushbottom` wyrównuje dolne krawędzie wszystkich stron, rozciągając
% w tym celu odstępy pionowe — przede wszystkim ten po nagłówku sekcji, bo jest
% najbardziej rozciągliwy. W kolumnie Kindle'a mieszczącej 24 wiersze brakujące
% dwa wiersze robią pod nagłówkiem **dziurę na jedną trzecią strony**. W logu jest
% to `Underfull \vbox while \output is active` — ostrzeżenie, nie błąd.
%
% Wyrównana dolna krawędź ma sens w książce oglądanej jako rozkładówka. Na
% czytniku widać jedną stronę naraz, więc nierówny dół jest niewidoczny, a dziura
% po nagłówku — bardzo.
\raggedbottom

% Pływaki w wąskiej kolumnie.
%
% Domyślne progi LaTeX-a zakładają stronę A4: rysunek zajmujący 70% wysokości nie
% ma prawa dzielić strony z tekstem, więc ląduje na osobnej stronie, a po sobie
% zostawia dziurę. Na ekranie Kindle'a 70% wysokości to zwykły wykres.
\renewcommand{\topfraction}{0.9}
\renewcommand{\bottomfraction}{0.9}
\renewcommand{\textfraction}{0.06}
\renewcommand{\floatpagefraction}{0.75}

% Listingi: łamanie długich wierszy kodu.
%
% Kod nie ma spacji tam, gdzie TeX chce łamać, więc w kolumnie szerokości
% Kindle'a wiersz po prostu **wychodzi poza stronę i zostaje ucięty** — bez
% błędu, tylko z ostrzeżeniem „Overfull hbox" w logu. Objaw jest podstępny,
% bo PDF wygląda poprawnie do chwili, gdy ktoś spróbuje przepisać kod.
%
% `breakanywhere` łamie także w środku identyfikatora — brzydkie, ale
% uczciwsze niż zgubiona połowa wiersza. `\small` w monospace daje ~62 znaki
% w wierszu zamiast 53, co dla większości kodu wystarcza bez łamania.
\usepackage{fvextra}
\fvset{breaklines=true, breakanywhere=true, fontsize=\small}
\DefineVerbatimEnvironment{Highlighting}{Verbatim}%
  {breaklines,breakanywhere,commandchars=\\\{\},fontsize=\small}
\DefineVerbatimEnvironment{verbatim}{Verbatim}%
  {breaklines,breakanywhere,fontsize=\small}

% Znak kontynuacji złamanego wiersza — bez niego nie widać, że wiersz kodu
% jest jednym wierszem, a nie dwoma.
\fvset{breaksymbolleft=\raisebox{0.6ex}{\tiny\ensuremath{\hookrightarrow}}}
